

%% ===================================================================
\section{연구가설}
\label{sec:hypotheses}
%% ===================================================================

$\DRTO$ 프레임워크의 이론적 타당성과 실증적 유효성을 검증하기 위해
\jrev{10개의 연구가설을 수립한다.} 가설은 시뮬레이션 검증\jrev{(H1--H3, 3건)},
회귀 검증(H4--H5, 2건), 강건성 검증(H6--H7, 2건),
이중채널 비선형 효과(H8, 1건), 전문가 평가(H9--H10, 2건)의 5개 범주로 구성된다.

\chg{각 가설은 서로 다른 이론적 근거에서 도출된다. H1에서 H3까지와 H8은 본 장에서 확립한 모델의 해석적 구조, 즉 시그모이드 바운딩의 부호 조건과 경계 보장 및 포화 성질에서 연역되고, H4와 H5는 입력 변수와 응답 사이의 비선형 관계를 완전 2차 다항식으로 근사하는 반응표면법~\citep{box1951experimental}에, H6과 H7은 몬테카를로 시뮬레이션의 독립 시드 설계와 계수 안정성 원리에, H9와 H10은 정량과 정성 방법을 결합하는 혼합연구방법의 삼각검증 원리~\citep{creswell2018designing}에 각각 근거를 둔다.}

\chg{각 가설에 대해서는 도출 배경과 핵심 명제, 검증 방향의 세 요소를 일관되게
제시한다. 도출 배경은 왜 이 가설이 필요한가를, 핵심 명제는 그 수학적이고 통계적인
정의를, 검증 방향은 어느 장에서 어떻게 검증되는가를 밝힌다. 시뮬레이션 검증에
해당하는 H1에서 H3까지는 모형의 기본 동작을, 회귀 검증인 H4와 H5는 변수 간 통계적
관계를, 강건성 검증인 H6과 H7은 결과의 재현성을, 이중채널의 H8은 본 연구의 핵심 신규
발견을, 전문가 평가인 H9와 H10은 외부 타당성을 각각 담당한다.
\secref{sec:ch4_hypothesis_summary}에서 10개 가설이 모두 채택(PASS)으로 종합되며, 그
정량 근거는 가설별 정의와 시뮬레이션 및 회귀 결과로 뒷받침된다.}


%% -------------------------------------------------------------------
%% -------------------------------------------------------------------

\chg{가설 H1(관측성 효과)은 관측성이 활성화된 조건(C1)에서 $\DRTO$가 정적 기준선(C0)보다 작거나 같다는 것이다.}
\[
  H_1:\;\; \DRTO_{\textbf{C1}} \leq \DRTO_{\textbf{C0}} = R_0
  \quad (\text{준수율 } 100\%)
\]

H1은 관측성의 효과 방향성을 검증한다.
수학적으로, C1 조건에서 $z_O \geq 0$이므로 $S(z_O) \geq 0$이 되어
$E_{O,\text{bnd}} = -R_0 \cdot S(z_O) \leq 0$이 항상 성립한다.
따라서 $\DRTO_{\text{C1}} = R_0 + E_{O,\text{bnd}} \leq R_0$가
모델 구조에 의해 해석적으로 보장된다.


\chg{가설 H2(시그모이드 경계 보장)는 모든 조건(C0에서 C3까지)에서 $\DRTO$가 물리적 경계 내에 위치한다는 것이다.}
\[
  H_2:\;\; 0 < \DRTO < R_{\max} = 4R_0 = 24.0\text{h}
  \quad (\text{준수율 } 100\%)
\]

H2는 개별 시그모이드 바운딩 모델의 핵심 설계 목표인
물리적 경계 보장을 검증한다. (명제 2.1)에 의해
이 경계는 $\kappa$에 무관하게 항상 성립한다.



\chg{가설 H3(조건 순서 보존)는 네 조건이 다음의 세 가지의 순서 관계를 유지한다는 것이다. $H_{3\text{-}a}$는
$\DRTO_{\textbf{C1}} \leq \DRTO_{\textbf{C0}}$(준수율 $100\%$), $H_{3\text{-}b}$는
$\DRTO_{\textbf{C1}} \leq \DRTO_{\textbf{C2}}$(준수율 $\geq 99\%$), $H_{3\text{-}c}$는
$\DRTO_{\textbf{C2}} \leq \DRTO_{\textbf{C3}}$(조건부, 백분위 $\geq$ P85.8)을 의미한다.}

H3는 조건 간 $\DRTO$의 크기 순서가 이론적 예측과 일치하는지를 검증한다.
H3-a는 H1과 동치이며 모델 구조에 의해 보장된다.
H3-b는 불확실성의 도입이 관측성의 하향 효과를 상쇄하고 초과함을 의미하며,
$N(3,1)$의 음수 표본(약 $0.135\%$)에서 $E_{U,\text{bnd}} < 0$이 되는 경우를
허용하기 위해 준수율 기준을 100\%가 아닌 99\% 이상으로 설정한다.
H3-c는 파레토 분포 특성이 가우시안 대비 높은 $\DRTO$를 산출함을
검증하나, 이 효과는 P85.8 CDF 교차점 이상의 극단 영역에서만 발현되므로
조건부 가설로 설정한다.


%% -------------------------------------------------------------------
%% -------------------------------------------------------------------

\chg{가설 H4(회귀 설명력)는 교호 작용을 포함한 2차 회귀 모델의 결정계수가 모든 조건에서 $R^2_{(2\mathrm{q})} \geq 0.95$를 달성한다는 것이다.}
\[
  \jrev{H_4}:\;\; R^2_{(2\mathrm{q})} \geq 0.95
  \quad \forall\; \textbf{C1}, \textbf{C2}, \textbf{C3}
\]

\erf{H5}{H4}는 $\DRTO$ 모델의 비선형 응답 곡면이 2차 다항식으로 충분히
근사될 수 있는지를 검증한다. 2차 회귀는 1차항($k_O$, $O$, $U$),
2차항($k_O^2$, $O^2$, $U^2$), 교호작용항($k_O \cdot O$, $k_O \cdot U$,
$O \cdot U$)을 포함하는 완전 2차 다항 모델이다.
\chg{이러한 모델 선택의 이론적 근거는 \citet{box1951experimental}이 정립한 반응표면법에 있다. 반응표면법은 입력 변수와 응답 사이의 미지의 비선형 관계를 국소적으로 곡률까지 포착하는 완전 2차 다항식으로 근사하는 접근이 타당함을 보였으며, 본 연구의 응답 곡면 또한 시그모이드 바운딩에서 비롯한 매끄러운 곡률 구조를 가지므로 이 근사가 적합하다.}


\chg{가설 H5(교호 작용 유의성)는 교호작용항(interaction terms)이 회귀 모델의 설명력에 유의하게 기여한다는 것이다.}
\[
  \jrev{H_5}:\;\; \Delta R^2_{(\text{교호작용항 추가})} > 0,\quad F\text{-검정에서 } p < 0.05
\]

H5는 관측성과 불확실성이 독립적으로 작용하는 것이 아니라,
두 요소의 곱(교호 작용)이 $\DRTO$에 추가적인 설명력을 제공함을 검증한다.
특히 $k_O \cdot O_{\text{raw}}$ 교호작용항은 C1에서 $R^2$를 결정적으로 향상시킨다.


%% -------------------------------------------------------------------
%% -------------------------------------------------------------------

\chg{가설 H6(비중첩 대역 시드 독립성)은 비중첩 대역(non-overlapping band) 시드에서 파생된 입력 변수 $k_O$, $O_{\text{raw}}$, $U_{\text{raw}}$ 간 상관계수의 \m{절댓값}이 0.02 미만이라는 것이다.}
\[
  \jrev{H_6}:\;\; |r(X_i, X_j)| < 0.02 \quad \forall\; i \neq j
\]

H6은 시뮬레이션의 입력 변수들이 통계적으로 독립임을 검증한다.
마스터 시드 42로부터 비중첩 대역 규칙에 따라 파생된 시드 간
상관계수가 무시 가능한 수준임을 확인한다.


\chg{가설 H7(다중 시드 강건성)은 다중 시드($s = 0, 1, \ldots, 9{,}999$) 시뮬레이션에서 회귀 계수가 안정적이라는 것이다.}
\[
  \jrev{H_7}:\;\; \mathrm{CV}(\beta_{k_O}) < 0.05
\]

H7은 seed = 42에서 도출된 결과가 특정 시드에 의한 우연적 산물이 아님을
검증한다. 10{,}000개의 독립 시드 각각에서 $n = 10{,}000$ 시뮬레이션을
수행하여 총 $10^8$회의 $\DRTO$ 계산을 실시하고,
회귀계수의 분포를 산출한다.


%% -------------------------------------------------------------------
%% -------------------------------------------------------------------

\chg{가설 H8(이중채널 비선형 감쇠)은 C3 꼬리 영역($\geq \text{P85.8}$)에서 관측성 가중치 $k_O$의 회귀 계수 크기가 핵심 영역($< \text{P85.8}$) 대비 감소한다(감쇠비 $< 1$)는 것이다.}
\[
  \jrev{H_8}:\;\; \frac{|\beta_{k_O}^{\text{Tail}}|}{|\beta_{k_O}^{\text{Core}}|} < 1
\]

H8은 본 연구에서 새롭게 발견한 이중 채널 비선형 상호작용(dual-channel nonlinear
interaction) 메커니즘을 형식화한 가설이다.
``이중 채널''이란 관측성 채널($z_O$)과 불확실성 채널($z_U$)의 두 경로를 의미하며,
두 채널이 시그모이드 함수의 비선형성을 통해 결합될 때
극단 영역에서 시그모이드 포화(saturation)로 인해 관측성의 한계 효과가
오히려 감소하는 현상을 지칭한다.
감쇠비의 실증값과 P85.8 분할 회귀 분석 상세 내용은 4.10(이중채널 감쇠 메커니즘)에서 다루며,
이는 극단적 불확실성 하에서 관측성 투자의 한계 효과가 체감함을 의미한다.
이러한 비선형 감쇠 메커니즘은 조직의 위험 관리 전략에 대한
중요한 정책적 시사점을 제공하며 5.3(실무적 기여)에서 상세히 논의한다.


%% -------------------------------------------------------------------
%% -------------------------------------------------------------------

\chg{아홉째 가설 H9(전문가 정량적 검증)는 BCM/DR 전문가 설문조사에서 전체 영역 평균이 5점 리커트 척도 기준 4.0 이상이라는 것이다.}
\[
  \jrev{H_{9}}:\;\; \bar{X}_{\text{expert}} \geq 4.0 \;/\; 5.0
\]


\chg{열째 가설 H10(전문가 정성적 합의)은 전문가 심층 인터뷰에서 도출된 핵심 주제에 대한 동의율이 $80\%$ 이상이라는 것이다. 각 주제 $j$에 대한 동의 여부를 이진 변수 $a_j \in \{0, 1\}$로 정의하면 $a_j = 1$은 동의를, $a_j = 0$은 비동의를 가리킨다.}
\[
  \jrev{H_{10}}:\;\; A_{\text{agree}} \geq 0.80, \quad
  A_{\text{agree}} = \frac{1}{N_{\text{topic}}}\sum_{j=1}^{N_{\text{topic}}} a_j
\]
\chg{여기서 $A_{\text{agree}}$는 전체 $N_{\text{topic}}$개 주제 중 전문가가 동의한 주제의 비율을 의미한다.}

H9와 H10은 $\DRTO$ 모델의 실무적 타당성을 현장 전문가의 판단을 통해
삼각검증(triangulation)하기 위한 가설로서, \chg{정량적 방법과 정성적 방법을 결합하여 단일 방법의 한계를 보완하는 혼합연구방법의 삼각검증 원리~\citep{creswell2018designing}에 이론적 근거를 둔다.}
정량적 설문(H9)은 4개 평가 영역, \chg{곧 모델 타당성과 파라미터 적정성,
실무 적용성, 산업적용 및 $\eta$ 활용에 대한} 27개 문항으로 구성되며,
정성적 인터뷰(H10)는 반구조화 방식으로 수행한다.
\chg{H9와 H10의 검증 결과 상세내용은 4.11(전문가 패널 검증)에서 다룬다.}


%% -------------------------------------------------------------------
%% -------------------------------------------------------------------

\p{[}표~\ref{tab:hypotheses-summary}\p{]}에 \jrev{10개} 가설의 정의, 판정 기준, 관련 조건을 종합하였다.
\chg{[표 2-16]에서 H7의 $10^8$은 다중 시드 시뮬레이션의 총 횟수를 가리킨다.}

\begin{table}[htbp]
\centering
\caption{\jrev{10개 연구가설 종합}}
\label{tab:hypotheses-summary}
\small
\begin{tabularx}{\textwidth}{c c X l c}
\toprule
\textbf{가설} & \textbf{범주} & \textbf{핵심 내용} & \textbf{판정 기준} & \textbf{조건} \\
\midrule
H1 & 시뮬레이션 & $\DRTO_{\text{C1}} \leq R_0$ (관측성 효과) & 준수율 100\% & C0, C1 \\
H2 & 시뮬레이션 & $0 < \DRTO < R_{\max} = 24.0$h (경계 보장) & 준수율 100\% & C0--C3 \\
H3 & 시뮬레이션 & C0--C3 순서 보존 (a/b/c) & 준수율 기준 & C0--C3 \\
\midrule
H4 & 회귀 & $R^2_{(2\mathrm{q})} \geq 0.95$ & C1/C2/C3 전체 & C1--C3 \\
H5 & 회귀 & 교호작용항 \m{유의성($\Delta R^2$)} & 기여도 검정 & C1--C3 \\
\midrule
H6 & 강건성 & 시드 간 $|r| < 0.02$ (독립성) & 상관계수 검정 & 전체 \\
H7 & 강건성 & 다중 시드 계수 \m{안정성($10^8$회)} & 계수 CV 기준 & 전체 \\
\midrule
H8 & 이중채널 & Tail/Core $k_O$ 감쇠비 $< 1$ & P85.8 분할 회귀 & C3 \\
\midrule
H9 & 전문가 & 설문 평균 $\geq 4.0/5.0$ & $t$-검정 & 전체 \\
H10 & 전문가 & $A_{\text{agree}} \geq 0.80$ (15개 하위 주제) & 합의 기준 & 전체 \\
\bottomrule
\end{tabularx}
\end{table}
